\section{第4章の解答}\label{解答:第4章}

\paragraph{4.1}
1つの小区間の幅は$\Delta x=\frac2n$,第$k$区間の右端は$x_k=\frac{2k}{n}$である.したがってリーマン和は
\begin{align*}
 \sum_{k=1}^n\left(\frac{2k}{n}\right)^2\frac2n
 &=\frac8{n^3}\sum_{k=1}^nk^2\\
 &=\frac8{n^3}\cdot\frac{n(n+1)(2n+1)}{6}\\
 &=\frac{4n(n+1)(2n+1)}{3n^3}
\end{align*}
$n\to\infty$で右辺は$\frac83$へ近づく.\,$x^2$は閉区間で連続なので,この極限が定積分の値であり,\,$\int_0^2x^2\dd x=\frac83$となる.

\paragraph{4.2}
区間の加法性から$\int_0^3f(x)\dd x=2+(-1)=1$である.次に積分の向きを直してから線形性を使うと,
\begin{align*}
 \int_3^0\{2f(x)-3g(x)\}\dd x
 &=-\left\{2\int_0^3f(x)\dd x-3\int_0^3g(x)\dd x\right\}\\
 &=-(2\cdot1-3\cdot4)\\
 &=10
\end{align*}
となる.
向きを逆にした負号は,被積分関数全体の積分に掛かる.

\paragraph{4.3}
$x-1$は$[-1,1]$で$0$以下,\,$[1,2]$で$0$以上である.下側の三角形は底辺と高さがともに$2$で面積$2$,上側の三角形は底辺と高さがともに$1$で面積$\frac12$となる.したがって定積分は$-2+\frac12=-\frac32$,幾何学的な面積は$2+\frac12=\frac52$である.積分で計算しても
\[
 \int_{-1}^{2}(x-1)\dd x
 =\left[\frac{x^2}{2}-x\right]_{-1}^{2}=-\frac32
\]
と一致する.面積を積分で求めるなら,負の区間の符号を反転して足す.

\paragraph{4.4}
いずれも区間内で連続であり,基本定理を適用できる.
\begin{align*}
 \text{(1)}\quad\int_{-1}^2(x^2+2x)\dd x
 &=\left[\frac{x^3}{3}+x^2\right]_{-1}^2\\
 &=\frac{20}{3}-\frac23\\
 &=6\\
 \text{(2)}\quad\int_0^{\frac\pi2}(\sin x+\cos x)\dd x
 &=[-\cos x+\sin x]_0^{\frac\pi2}\\
 &=1-(-1)\\
 &=2\\
 \text{(3)}\quad\int_1^{\e}\left(\frac2x+x\right)\dd x
 &=\left[2\log x+\frac{x^2}{2}\right]_1^{\e}\\
 &=\left(2+\frac{\e^2}{2}\right)-\frac12\\
 &=\frac{\e^2+3}{2}
\end{align*}
(3)の区間は正の側にあるので,対数を$\log x$と書ける.

\paragraph{4.5}
基本定理から$A'(x)=(x-2)(x+1)$である.\,$[-2,-1]$で増加,\,$[-1,2]$で減少,\,$[2,3]$で増加するため,閉区間での最大・最小の候補は両端と$x=-1,2$である.比較するために
\[
 A(x)=\left[\frac{t^3}{3}-\frac{t^2}{2}-2t\right]_1^x
 =\frac{x^3}{3}-\frac{x^2}{2}-2x+\frac{13}{6}
\]
と求めると,
\[
 A(-2)=\frac32,\quad A(-1)=\frac{10}{3},\quad
 A(2)=-\frac76,\quad A(3)=\frac23.
\]
よって最大値は$x=-1$で$\frac{10}{3}$,最小値は$x=2$で$-\frac76$である.導関数から増減を先に調べたので,すべての点で値を比較する必要はない.

\paragraph{4.6}
上端は$\sin x$,下端は$\cos x$である.それぞれの微分まで掛けると
\[
 G'(x)=\sqrt{1+\sin^2x}\cos x
       -\sqrt{1+\cos^2x}(-\sin x)
 =\sqrt{1+\sin^2x}\cos x+\sqrt{1+\cos^2x}\sin x.
\]
下端を微分したときの負号と,端点の公式にある負号が掛かって正になる.上下端の大小が入れ替わっても同じ式が成り立つ.

\paragraph{4.7}
$x^2\ge0$だから$\sqrt{1+x^2}\ge1$である.また両辺が非負であり,
\[
 \left(1+\frac{x^2}{2}\right)^2-(1+x^2)=\frac{x^4}{4}\ge0
\]
なので,平方の大小から$\sqrt{1+x^2}\le1+\frac{x^2}{2}$を得る.\,$[0,1]$上でこの不等式を積分すると
\[
 1=\int_0^1 1\dd x
 \le\int_0^1\sqrt{1+x^2}\dd x
 \le\int_0^1\left(1+\frac{x^2}{2}\right)\dd x
 =1+\frac16=\frac76.
\]
近似の式を単に代入したのではなく,区間全体の不等式を確かめてから積分している.

\paragraph{4.8}
(1) $x=0$を代入すると$f(0)=3$である.右辺は連続な$f$の積分なので微分可能であり,等式から$f$も微分可能となる.微分すれば$f'=-f$だから
\[
 (\e^xf(x))'=\e^x\{f(x)+f'(x)\}=0.
\]
したがって$\e^xf(x)$は定数で,初期値から$f(x)=3\e^{-x}$となる.元の式にも
\[
 3-\int_0^x3\e^{-t}\dd t
 =3-[-3\e^{-t}]_0^x=3\e^{-x}
\]
と戻るので,これが求める関数である.

(2) 積分$K=\int_0^1g(t)\dd t$は$x$に依存しない定数である.条件は$g(x)=x+K$と書ける.しかし,この式を$0$から$1$まで積分すると
\[
 K=\int_0^1(x+K)\dd x=\frac12+K
\]
となり矛盾する.したがって条件を満たす連続関数は存在しない.微分だけなら$g'=1$を得るが,それだけでは失われた定数の条件を確かめられない.
